% Auxiliary coordinate statements for polynomial PEPS compression.
% These statements do not assert the sampled-contraction identity or Theorem 5.2.

\subsection{Physical bases and party grouping}

Let $P$ be a finite party set and let $p_1,\ldots,p_m$ be an
ordering containing each party exactly once. For a register list $a$,
write $M(a)$ for its right-associated tensor memory, including a terminal
factor $\C$, and $a_p$ for the registers owned by $p$, in their original
order. The canonical isometry $G_a$ successively partitions off
$a_{p_1},\ldots,a_{p_m}$, identifying $M(a)$ with
$M(a_{p_1})\otimes\cdots\otimes M(a_{p_m})\otimes\C$.

\begin{definition}[Party tensor bases]
    \label{def:peps_party_tensor_basis}
    \lean{TNLean.PEPS.PairEffect.partyListBasis,
        TNLean.PEPS.PairEffect.partyLabelledBasis}
    \leanok
    \uses{def:peps_party_layout}
    For each $p$, let $b_p$ be an orthonormal basis of $M(a_p)$,
    indexed by a finite set $I_p$. The ordered party basis has column
    $b_{p_1}(x_{p_1})\otimes\cdots\otimes b_{p_m}(x_{p_m})\otimes1$,
    indexed by assignments $x_p\in I_p$. For an empty ordering,
    its single column is the scalar $1$.
\end{definition}

\begin{definition}[Party bases on the original register order]
    \label{def:peps_party_grouped_basis}
    \lean{TNLean.PEPS.PairEffect.partyGroupedBasis}
    \leanok
    \uses{def:peps_party_tensor_basis}
    Transport the ordered party basis by $G_a^{-1}$. Its column on the
    original register memory is
    \begin{align}
        B^a_x=G_a^{-1}
        \bigl(b_{p_1}(x_{p_1})\otimes\cdots
        \otimes b_{p_m}(x_{p_m})\otimes1\bigr).
    \end{align}
\end{definition}

\begin{theorem}[Singleton physical memory of each party]
    \label{thm:peps_singleton_physical_memory}
    \lean{TNLean.PEPS.PairEffect.atParty_familyPhysicalLayout_of_mem}
    \leanok
    Let an ordered list of parties have no repetitions. For any party
    $p$ in this list, selecting its physical registers gives exactly
    its one register $H_p$, with memory $H_p\otimes\C$.
\end{theorem}

\begin{proof}
    \leanok
    Selecting a party filters the given list by equality with $p$.
    Since $p$ occurs exactly once, the filtered register list is the
    singleton containing its physical register.
\end{proof}

\begin{definition}[Physical basis at an individual party]
    \label{def:peps_singleton_physical_basis}
    \lean{TNLean.PEPS.PairEffect.familyPhysicalPartyBasis}
    \leanok
    \uses{thm:peps_singleton_physical_memory}
    On the physical memory $H_p\otimes\C$ of party $p$, take the
    orthonormal basis with columns $b_p(j)=e_{p,j}\otimes1$.
    This is the standard basis transported by the inverse of the
    right scalar-unit isometry.
\end{definition}

\begin{theorem}[Partitioning a single selected register]
    \label{thm:peps_selected_head_partition}
    \lean{TNLean.PEPS.PairEffect.Layout.partitionIso_selectedHead_of_tail_excluded_heq}
    \leanok
    Let the first register have owner $p$ and Hilbert space $H$, and
    let every remaining register have an owner different from $p$.
    Write $K$ for the memory of the remaining registers. Under the
    canonical partition into the registers of $p$ and their complement,
    every $e\in H$ and $y\in K$ satisfy
    \begin{align}
        e\otimes y\longmapsto(e\otimes1)\otimes y.
    \end{align}
\end{theorem}

\begin{proof}
    \leanok
    Partition the remaining registers first. Their selected memory is
    $\C$, so $y$ is sent to $1\otimes y$. Adjoining the first register
    and applying the associativity isometry gives the stated formula.
\end{proof}

\begin{theorem}[Grouping prescribed physical columns in arbitrary local coordinates]
    \label{thm:peps_basis_physical_columns_local_coordinates}
    \lean{TNLean.PEPS.PairEffect.groupByPartyIso_familyPhysicalListBasis_eq_partyListBasis}
    \leanok
    \uses{def:peps_family_physical_output,def:peps_party_tensor_basis}
    Let $P$ be an arbitrary set of parties, let $d_p$ be arbitrary
    nonnegative integers, and let $\mathbf p=(p_0,\ldots,p_{m-1})$
    be an ordered list without repetitions. Put
    $a=a_{\mathbf d}(\mathbf p)$, with register $\C^{d_{p_j}}$
    at the party in position $j$. For each party $p$, let $b_p$
    be an orthonormal basis of the complete memory $M(a_p)$,
    indexed by a finite set $I_p$. Let
    $x\in\prod_{p\in P}\{0,\ldots,d_p-1\}$ and
    $y\in\prod_{p\in P} I_p$.

    Suppose that, at every party in $\mathbf p$, the canonical
    singleton-memory identification sends $b_p(y_p)$ to
    $e_{x_p}\otimes1$. Then grouping the prescribed physical tensor
    column gives the ordered column of these local bases:
    \begin{align}
        G_{\mathbf p,a}\,
        b_{\mathbf d,\mathbf p}\bigl((x_{p_j})_{j=0}^{m-1}\bigr)
        = b_{a,\mathbf p}\bigl((y_{p_j})_{j=0}^{m-1}\bigr).
    \end{align}
    Here $b_{\mathbf d,\mathbf p}$ is the ordered Euclidean tensor
    basis, $G_{\mathbf p,a}$ successively separates the complete memories
    in the listed order, and $b_{a,\mathbf p}$ is the ordered tensor
    basis formed from $b_{p_0},\ldots,b_{p_{m-1}}$.
    The local index sets need not be the physical index sets.
    No coverage of all of $P$ is required: the list covers the owners
    of $a$ by construction. The empty list and zero local dimensions
    are included, with quantification over the indicated assignments.
    Both ordered memories retain their terminal scalar factor $\C$.
\end{theorem}

\begin{proof}\leanok
    \uses{thm:peps_selected_head_partition,def:peps_party_tensor_basis}
    Induct on the length of the ordered list. For the empty list,
    both columns are $1\in\C$. At a nonempty list, its first
    party does not occur in the tail. The selected-register partition
    therefore separates its physical column as $e_{x_p}\otimes1$.
    Apply the induction hypothesis to the tail. The prescribed physical
    column thus groups as the ordered tensor product of its singleton
    columns. By the given identifications, these are precisely
    $b_{p_j}(y_{p_j})$ at every listed party. The ordered tensor basis
    has the same column,
    \begin{align}
        b_{p_0}(y_{p_0})\otimes\cdots\otimes
        b_{p_{m-1}}(y_{p_{m-1}})\otimes1.
    \end{align}
\end{proof}

\begin{theorem}[Compatibility of physical and party coordinates]
    \label{thm:peps_family_physical_party_basis}
    \lean{TNLean.PEPS.PairEffect.familyLabelledPhysicalBasis_eq_partyGroupedBasis}
    \leanok
    Let the physical register list $a$ and the singleton bases $b_p$
    be as above. For every physical assignment $x$,
    \begin{align}
        G_a E_x
         & =(e_{p_1,x_{p_1}}\otimes1)\otimes\cdots
        \otimes(e_{p_m,x_{p_m}}\otimes1)\otimes1.
    \end{align}
    Thus the labelled physical basis equals the party basis on the
    original register order. No positivity of $d_p$ is required.
    If some $d_p=0$, the assignment set and both bases are empty.
    If $P$ is empty, both tensor products are the scalar $1$.
\end{theorem}

\begin{proof}
    \leanok
    \uses{def:peps_party_grouped_basis,def:peps_family_physical_output,
        def:peps_singleton_physical_basis,thm:peps_selected_head_partition}
    First consider any ordered list without repetitions, retaining
    only its physical registers. Induct on this list. For the empty
    list, both memories are $\C$ with basis vector $1$. In the
    induction step, no remaining register is owned by the first party.
    Theorem~\ref{thm:peps_selected_head_partition} separates its basis
    vector as $e_{p,j}\otimes1$, leaving the ordered physical basis
    vector on the tail. Apply the induction hypothesis to that tail.
    When the list exhausts $P$, transporting the resulting equality
    by $G_a^{-1}$ gives equality of the two bases.
\end{proof}

\begin{definition}[Joint physical and private basis of one party]
    \label{def:peps_party_local_output_basis}
    \lean{TNLean.PEPS.PairEffect.partyLocalOutputBasis}
    \leanok
    For register lists $a,c$, choose orthonormal bases $b_p$ of
    $M(a_p)$ and $f_p$ of $M(c_p)$, indexed respectively by finite
    sets $X_p$ and $F_p$. Let $J_{u,v}:M(u\mathbin{+\!+}v)
        \longrightarrow M(u)\otimes M(v)$ be the canonical concatenation
    isometry. The joint basis at party $p$ has column
    $J_{a_p,c_p}^{-1}(b_p(x_p)\otimes f_p(y_p))$, under the
    identification $(a\mathbin{+\!+}c)_p=a_p\mathbin{+\!+}c_p$.
\end{definition}

\begin{theorem}[Concatenating physical and private party bases]
    \label{thm:peps_party_grouped_basis_append}
    \lean{TNLean.PEPS.PairEffect.partyGroupedBasis_append}
    \leanok
    With the joint local bases just defined, let $B^{a+c}$ denote
    their party basis on $M(a\mathbin{+\!+}c)$, and let $B^a,B^c$
    be the separately transported party bases. Then for all
    assignments $x_p\in X_p$ and $y_p\in F_p$,
    \begin{align}
        B^{a+c}_{(x,y)}=J_{a,c}^{-1}(B^a_x\otimes B^c_y).
    \end{align}
    The index $(x,y)$ denotes the assignment $p\mapsto(x_p,y_p)$;
    equivalently, this is equality of the two bases after the canonical
    separation of the two assignment families.
\end{theorem}

\begin{proof}
    \leanok
    \uses{def:peps_party_grouped_basis,def:peps_party_local_output_basis}
    Induct on the ordered party list. Partitioning a concatenation
    first partitions each list and then regroups its four factors,
    joining the two selected memories and the two complementary
    memories. The selected vector is the stated joint local column.
    Apply the induction hypothesis to the two complementary lists,
    and transport by $G_{a+c}^{-1}$. For the empty list all memories
    are scalar and the assertion reduces to multiplication in $\C$.
\end{proof}

These coordinate identities apply to the local contractions in
\cite[proof of Theorem~5.2]{OpenAI2026PolynomialPEPS},
\texttt{04-compression.tex}, lines 137--151 and 565--588.
